\clearpage
\section{Appendix: Scriptural Date and Location}
\label{historical-appendix}
% Generated. Do not edit.
% Deterministic projection of research/chronology.toml.
% schema: 5
% document: liturgy/roman-rite/postconciliar/roman-missal-third-edition-en-us-2011/propers/temporal/pc-s55-twenty-ninth-sunday-in-ordinary-time-year-a
% generated-by: tools/tpt proper-chronology annotations
\newcommand{\chronologyannotationclaim}[6]{#6}
\newcommand{\chronologyannotationreach}[3]{}
\newcommand{\chronologyannotationgroup}[3]{#3}
\newcommand{\chronologyannotation}[1]{%
  \ifcsname triptychchronologyannotation@#1\endcsname
    \csname triptychchronologyannotation@#1\endcsname
  \else
    \PackageError{triptych}{No chronology annotation for #1}{}%
  \fi}
% appointment-dependency: guidance/liturgy/postconciliar-propers-registry.md sha256=ffe8a6d8b7d79c652bdc49b3d5f6d4dbb32e3506abc7ab0310e745c2b7306ad4
% appointment-dependency: src/gpt/liturgy/roman-rite/postconciliar/roman-missal-third-edition-en-us-2011/propers/registry/README.md sha256=454b7262c8e30f246a46629e6da476242a7166e6cb38a2370eeda52c72c5554a
% appointment-dependency: src/gpt/liturgy/roman-rite/postconciliar/roman-missal-third-edition-en-us-2011/propers/registry/formula-dispositions.md sha256=281e5dd2484b732176f3d4136a070527dbbdb6abbcfc0adfda49e08cd90d206d
% appointment-dependency: src/gpt/liturgy/roman-rite/postconciliar/roman-missal-third-edition-en-us-2011/propers/temporal/pc-s55-twenty-ninth-sunday-in-ordinary-time-year-a/instance/manifest.md sha256=a33c95bbbb2f9b954e8f879cdcb159ed91b18220a86e6917c9a6092c1e6d32fd
% appointment-dependency: src/gpt/liturgy/roman-rite/postconciliar/roman-missal-third-edition-en-us-2011/propers/temporal/pc-s55-twenty-ninth-sunday-in-ordinary-time-year-a/propers/verified.md sha256=4e5b677a3ceeffc388e41c1d9e2e31b690200ab9ba63365e88e975229db7de2a
% appointment-dependency: src/gpt/liturgy/roman-rite/postconciliar/roman-missal-third-edition-en-us-2011/propers/temporal/pc-s55-twenty-ninth-sunday-in-ordinary-time-year-a/research/chronology-inputs.toml sha256=0b689b16340810b8d2a46dcf681626b0b7ba23be25e11a5f8fcf033b1d33dce4
% appointment-dependency: src/gpt/liturgy/roman-rite/postconciliar/roman-missal-third-edition-en-us-2011/propers/temporal/shared/ordinary-time/weeks/29/propers/verified.md sha256=1e6ac1e5369a1ab134cc2969d1caf70f7dead1c18610987f6b89d80c11d089b1
% appointment-note: entrance: Cf. Psalm 16:6,8; adaptation; vulgate numbering; canonical verse queries, not appointed wording.
% appointment-note: collect: Composed Missal oration without an appointed biblical citation; theological affinities do not create appointments.
% appointment-note: first-reading: Isaiah 45:1,4-6; appointed; vulgate numbering; canonical verse queries, not appointed wording.
% appointment-note: psalm: Psalm 96:1,3,4-5,7-8,9,10a,10c; appointed; hebrew numbering; whole-verse query envelope, not subverse chronology or appointed wording.
% appointment-note: psalm: Psalm 96:7b; response; hebrew numbering; whole-verse query envelope, not subverse chronology or appointed wording.
% appointment-note: second-reading: 1 Thessalonians 1:1-5b; appointed; vulgate numbering; whole-verse query envelope, not subverse chronology or appointed wording.
% appointment-note: acclamation: Philippians 2:15d,16a; appointed; vulgate numbering; whole-verse query envelope, not subverse chronology or appointed wording.
% appointment-note: gospel: Matthew 22:15-21; appointed; vulgate numbering; canonical verse queries, not appointed wording.
% appointment-note: offerings: Composed Missal oration without an appointed biblical citation; theological affinities do not create appointments.
% appointment-note: communion-psalm: Cf. Psalm 32:18-19; adaptation; vulgate numbering; canonical verse queries, not appointed wording.
% appointment-note: communion-mark: Mark 10:45; identified-basis; vulgate numbering; canonical verse queries, not appointed wording.
% appointment-note: after-communion: Composed Missal oration without an appointed biblical citation; theological affinities do not create appointments.
% entrance: Entrance antiphon
\expandafter\def\csname triptychchronologyannotation@entrance\endcsname{%
  \chronologyannotationgroup{historical-setting}{research-pending}{\textbf{Historical setting} -- Occasion date unresolved.}%
  \space \chronologyannotationgroup{composition}{preferred}{\textbf{Composition}: \chronologyannotationreach{Ps.16.6}{Ps}{true}\chronologyannotationreach{Ps.16.8}{Ps}{true}\chronologyannotationclaim{critical.psalms.latest-composition-boundary}{composition}{catholic-critical-v1}{preferred}{before the Maccabean period, around 165 B.C.; no individual psalm can be dated securely}{Before c. 165 B.C.}}%
}
% first-reading: First reading
\expandafter\def\csname triptychchronologyannotation@first-reading\endcsname{%
  \chronologyannotationgroup{traditional-attribution}{preferred}{\textbf{Traditional attribution}: \chronologyannotationreach{Is.45.1}{Is}{true}\chronologyannotationreach{Is.45.4}{Is}{true}\chronologyannotationreach{Is.45.5}{Is}{true}\chronologyannotationreach{Is.45.6}{Is}{true}\chronologyannotationclaim{israel.prophets.isaias-traditional-era}{traditional-attribution}{catholic-traditional-v1}{preferred}{from the closing year of Ozias, King of Juda (740 B.C.); no prophecies appear to be later than 701 B.C.}{Isaias (ministry in Souvay's traditional account), B.C. 740--701}.}%
  \space \chronologyannotationgroup{prophecy-given}{preferred}{\textbf{Prophecy given}: \chronologyannotationreach{Is.45.1}{Is.40 Is.41 Is.42 Is.43 Is.44 Is.45 Is.46 Is.47 Is.48 Is.49 Is.50 Is.51 Is.52 Is.53 Is.54 Is.55}{true}\chronologyannotationreach{Is.45.4}{Is.40 Is.41 Is.42 Is.43 Is.44 Is.45 Is.46 Is.47 Is.48 Is.49 Is.50 Is.51 Is.52 Is.53 Is.54 Is.55}{true}\chronologyannotationreach{Is.45.5}{Is.40 Is.41 Is.42 Is.43 Is.44 Is.45 Is.46 Is.47 Is.48 Is.49 Is.50 Is.51 Is.52 Is.53 Is.54 Is.55}{true}\chronologyannotationreach{Is.45.6}{Is.40 Is.41 Is.42 Is.43 Is.44 Is.45 Is.46 Is.47 Is.48 Is.49 Is.50 Is.51 Is.52 Is.53 Is.54 Is.55}{true}\chronologyannotationclaim{israel.exile.second-isaias-oracles}{prophecy-given}{catholic-critical-v1}{preferred}{toward the end of the Babylonian exile}{The prophetic activity behind Isaiah 40-55 in the NABRE introduction, toward the end of the Babylonian exile}.}%
}
% psalm: Responsorial Psalm
\expandafter\def\csname triptychchronologyannotation@psalm\endcsname{%
  \chronologyannotationgroup{superscription-setting}{preferred}{\textbf{Superscription setting}: \chronologyannotationreach{Ps.95.1}{Ps.95}{true}\chronologyannotationreach{Ps.95.10}{Ps.95}{true}\chronologyannotationreach{Ps.95.3}{Ps.95}{true}\chronologyannotationreach{Ps.95.4}{Ps.95}{true}\chronologyannotationreach{Ps.95.5}{Ps.95}{true}\chronologyannotationreach{Ps.95.7}{Ps.95}{true}\chronologyannotationreach{Ps.95.8}{Ps.95}{true}\chronologyannotationreach{Ps.95.9}{Ps.95}{true}\chronologyannotationclaim{israel.restoration.temple-raised-by-zorobabel}{superscription-setting}{catholic-traditional-v1}{preferred}{(537 B.C.)}{The raising of the second Temple by Zorobabel, B.C. 537}.}%
  \space \chronologyannotationgroup{composition}{preferred}{\textbf{Composition}: \chronologyannotationreach{Ps.95.1}{Ps}{true}\chronologyannotationreach{Ps.95.10}{Ps}{true}\chronologyannotationreach{Ps.95.3}{Ps}{true}\chronologyannotationreach{Ps.95.4}{Ps}{true}\chronologyannotationreach{Ps.95.5}{Ps}{true}\chronologyannotationreach{Ps.95.7}{Ps}{true}\chronologyannotationreach{Ps.95.8}{Ps}{true}\chronologyannotationreach{Ps.95.9}{Ps}{true}\chronologyannotationclaim{critical.psalms.latest-composition-boundary}{composition}{catholic-critical-v1}{preferred}{before the Maccabean period, around 165 B.C.; no individual psalm can be dated securely}{Before c. 165 B.C.}}%
  \space \chronologyannotationgroup{prophetic-referent}{disputed}{\textbf{Prophetic referent} -- disputed: \chronologyannotationreach{Ps.95.1}{Ps.67 Ps.95 Ps.96 Ps.97}{true}\chronologyannotationreach{Ps.95.10}{Ps.67 Ps.95 Ps.96 Ps.97}{true}\chronologyannotationreach{Ps.95.3}{Ps.67 Ps.95 Ps.96 Ps.97}{true}\chronologyannotationreach{Ps.95.4}{Ps.67 Ps.95 Ps.96 Ps.97}{true}\chronologyannotationreach{Ps.95.5}{Ps.67 Ps.95 Ps.96 Ps.97}{true}\chronologyannotationreach{Ps.95.7}{Ps.67 Ps.95 Ps.96 Ps.97}{true}\chronologyannotationreach{Ps.95.8}{Ps.67 Ps.95 Ps.96 Ps.97}{true}\chronologyannotationreach{Ps.95.9}{Ps.67 Ps.95 Ps.96 Ps.97}{true}\chronologyannotationclaim{life-of-christ.nativity}{prophetic-referent}{catholic-traditional-v1}{disputed}{the year of Rome 750 which he styles 3 B.C.}{The Nativity of Our Lord, B.C. 3}; \chronologyannotationreach{Ps.95.1}{Ps.67 Ps.95 Ps.96 Ps.97}{true}\chronologyannotationreach{Ps.95.10}{Ps.67 Ps.95 Ps.96 Ps.97}{true}\chronologyannotationreach{Ps.95.3}{Ps.67 Ps.95 Ps.96 Ps.97}{true}\chronologyannotationreach{Ps.95.4}{Ps.67 Ps.95 Ps.96 Ps.97}{true}\chronologyannotationreach{Ps.95.5}{Ps.67 Ps.95 Ps.96 Ps.97}{true}\chronologyannotationreach{Ps.95.7}{Ps.67 Ps.95 Ps.96 Ps.97}{true}\chronologyannotationreach{Ps.95.8}{Ps.67 Ps.95 Ps.96 Ps.97}{true}\chronologyannotationreach{Ps.95.9}{Ps.67 Ps.95 Ps.96 Ps.97}{true}\chronologyannotationclaim{life-of-christ.nativity}{prophetic-referent}{catholic-traditional-v1}{disputed}{probably the year 7 B.C.}{The Nativity of Our Lord, c. B.C. 7}.}%
}
% second-reading: Second reading
\expandafter\def\csname triptychchronologyannotation@second-reading\endcsname{%
  \chronologyannotationgroup{composition}{disputed}{\textbf{Composition} -- disputed: \chronologyannotationreach{1Thess.1.1}{1Thess}{true}\chronologyannotationreach{1Thess.1.2}{1Thess}{true}\chronologyannotationreach{1Thess.1.3}{1Thess}{true}\chronologyannotationreach{1Thess.1.4}{1Thess}{true}\chronologyannotationreach{1Thess.1.5}{1Thess}{true}\chronologyannotationclaim{composition.first-epistle-to-the-thessalonians}{composition}{catholic-traditional-v1}{disputed}{(1 and 2 Thessalonians), 52}{A.D. 52}; \chronologyannotationreach{1Thess.1.1}{1Thess}{true}\chronologyannotationreach{1Thess.1.2}{1Thess}{true}\chronologyannotationreach{1Thess.1.3}{1Thess}{true}\chronologyannotationreach{1Thess.1.4}{1Thess}{true}\chronologyannotationreach{1Thess.1.5}{1Thess}{true}\chronologyannotationclaim{composition.first-epistle-to-the-thessalonians}{composition}{catholic-traditional-v1}{disputed}{in the year 53 or 54 according to the commonly received scheme of Pauline chronology}{A.D. 53--54}.}%
}
% acclamation: Gospel acclamation
\expandafter\def\csname triptychchronologyannotation@acclamation\endcsname{%
  \chronologyannotationgroup{composition}{disputed}{\textbf{Composition} -- disputed: \chronologyannotationreach{Phil.2.15}{Phil}{true}\chronologyannotationreach{Phil.2.16}{Phil}{true}\chronologyannotationclaim{composition.epistle-to-the-philippians}{composition}{catholic-traditional-v1}{disputed}{(Philemon; Colossians; Ephesians; Philippians), 61}{A.D. 61}; \chronologyannotationreach{Phil.2.15}{Phil}{true}\chronologyannotationreach{Phil.2.16}{Phil}{true}\chronologyannotationclaim{composition.epistle-to-the-philippians}{composition}{catholic-traditional-v1}{disputed}{at Rome (A.D. 62-64)}{A.D. 62--64}.}%
}
% gospel: Gospel
\expandafter\def\csname triptychchronologyannotation@gospel\endcsname{%
  \chronologyannotationgroup{narrated-event}{research-pending}{\textbf{Event} -- No narrated-event date in the chronology corpus.}%
  \space \chronologyannotationgroup{composition}{disputed}{\textbf{Composition} -- disputed: \chronologyannotationreach{Matt.22.15}{Matt}{true}\chronologyannotationreach{Matt.22.16}{Matt}{true}\chronologyannotationreach{Matt.22.17}{Matt}{true}\chronologyannotationreach{Matt.22.18}{Matt}{true}\chronologyannotationreach{Matt.22.19}{Matt}{true}\chronologyannotationreach{Matt.22.20}{Matt}{true}\chronologyannotationreach{Matt.22.21}{Matt}{true}\chronologyannotationclaim{composition.gospel-of-matthew}{composition}{catholic-traditional-v1}{disputed}{about A.D. 38-45}{c. A.D. 38--45}; \chronologyannotationreach{Matt.22.15}{Matt}{true}\chronologyannotationreach{Matt.22.16}{Matt}{true}\chronologyannotationreach{Matt.22.17}{Matt}{true}\chronologyannotationreach{Matt.22.18}{Matt}{true}\chronologyannotationreach{Matt.22.19}{Matt}{true}\chronologyannotationreach{Matt.22.20}{Matt}{true}\chronologyannotationreach{Matt.22.21}{Matt}{true}\chronologyannotationclaim{composition.gospel-of-matthew}{composition}{catholic-traditional-v1}{disputed}{about the year 40-42}{c. A.D. 40--42}; \chronologyannotationreach{Matt.22.15}{Matt}{true}\chronologyannotationreach{Matt.22.16}{Matt}{true}\chronologyannotationreach{Matt.22.17}{Matt}{true}\chronologyannotationreach{Matt.22.18}{Matt}{true}\chronologyannotationreach{Matt.22.19}{Matt}{true}\chronologyannotationreach{Matt.22.20}{Matt}{true}\chronologyannotationreach{Matt.22.21}{Matt}{true}\chronologyannotationclaim{composition.gospel-of-matthew}{composition}{catholic-traditional-v1}{disputed}{the years 40-45}{A.D. 40--45}; \chronologyannotationreach{Matt.22.15}{Matt}{true}\chronologyannotationreach{Matt.22.16}{Matt}{true}\chronologyannotationreach{Matt.22.17}{Matt}{true}\chronologyannotationreach{Matt.22.18}{Matt}{true}\chronologyannotationreach{Matt.22.19}{Matt}{true}\chronologyannotationreach{Matt.22.20}{Matt}{true}\chronologyannotationreach{Matt.22.21}{Matt}{true}\chronologyannotationclaim{composition.gospel-of-matthew}{composition}{catholic-traditional-v1}{disputed}{about the year 60-68}{c. A.D. 60--68}; \chronologyannotationreach{Matt.22.15}{Matt}{true}\chronologyannotationreach{Matt.22.16}{Matt}{true}\chronologyannotationreach{Matt.22.17}{Matt}{true}\chronologyannotationreach{Matt.22.18}{Matt}{true}\chronologyannotationreach{Matt.22.19}{Matt}{true}\chronologyannotationreach{Matt.22.20}{Matt}{true}\chronologyannotationreach{Matt.22.21}{Matt}{true}\chronologyannotationclaim{composition.gospel-of-matthew}{composition}{catholic-traditional-v1}{disputed}{about the years 64-67}{c. A.D. 64--67}; \chronologyannotationreach{Matt.22.15}{Matt}{true}\chronologyannotationreach{Matt.22.16}{Matt}{true}\chronologyannotationreach{Matt.22.17}{Matt}{true}\chronologyannotationreach{Matt.22.18}{Matt}{true}\chronologyannotationreach{Matt.22.19}{Matt}{true}\chronologyannotationreach{Matt.22.20}{Matt}{true}\chronologyannotationreach{Matt.22.21}{Matt}{true}\chronologyannotationclaim{composition.gospel-of-matthew}{composition}{catholic-traditional-v1}{disputed}{about the year 50}{c. A.D. 50}.}%
}
% communion-psalm: Communion: psalm option
\expandafter\def\csname triptychchronologyannotation@communion-psalm\endcsname{%
  \chronologyannotationgroup{traditional-attribution}{disputed}{\textbf{Traditional attribution} -- disputed: \chronologyannotationreach{Ps.32.18}{Ps.32}{true}\chronologyannotationreach{Ps.32.19}{Ps.32}{true}\chronologyannotationclaim{israel.monarchy.david-traditional-era}{traditional-attribution}{catholic-traditional-v1}{disputed}{reigned from 1055 to 1015 B.C.}{David (reign in the usual chronology), B.C. 1055--1015}.}%
  \space \chronologyannotationgroup{historical-setting}{research-pending}{\textbf{Historical setting} -- Occasion date unresolved.}%
  \space \chronologyannotationgroup{composition}{preferred}{\textbf{Composition}: \chronologyannotationreach{Ps.32.18}{Ps}{true}\chronologyannotationreach{Ps.32.19}{Ps}{true}\chronologyannotationclaim{critical.psalms.latest-composition-boundary}{composition}{catholic-critical-v1}{preferred}{before the Maccabean period, around 165 B.C.; no individual psalm can be dated securely}{Before c. 165 B.C.}}%
}
% communion-mark: Communion: Mark option
\expandafter\def\csname triptychchronologyannotation@communion-mark\endcsname{%
  \chronologyannotationgroup{composition}{disputed}{\textbf{Composition} -- disputed: \chronologyannotationreach{Mark.10.45}{Mark}{true}\chronologyannotationclaim{composition.gospel-of-mark}{composition}{catholic-traditional-v1}{disputed}{the tenth or twelfth year after the Ascension (A.D. 38-40)}{A.D. 38--40}; \chronologyannotationreach{Mark.10.45}{Mark}{true}\chronologyannotationclaim{composition.gospel-of-mark}{composition}{catholic-traditional-v1}{disputed}{A.D. 40}{A.D. 40}; \chronologyannotationreach{Mark.10.45}{Mark}{true}\chronologyannotationclaim{composition.gospel-of-mark}{composition}{catholic-traditional-v1}{disputed}{the third year of Claudius (A.D. 43)}{A.D. 43}; \chronologyannotationreach{Mark.10.45}{Mark}{true}\chronologyannotationclaim{composition.gospel-of-mark}{composition}{catholic-traditional-v1}{disputed}{somewhere between A.D. 50 and 67}{A.D. 50--67}; \chronologyannotationreach{Mark.10.45}{Mark}{true}\chronologyannotationclaim{composition.gospel-of-mark}{composition}{catholic-traditional-v1}{disputed}{between 60 and 70}{A.D. 60--70}; \chronologyannotationreach{Mark.10.45}{Mark}{true}\chronologyannotationclaim{composition.gospel-of-mark}{composition}{catholic-traditional-v1}{disputed}{after Matthew, and not as late as the destruction of Jerusalem}{After Matthew, and not as late as the destruction of Jerusalem}.}%
}
% presentation: orientation-v1
\newcommand{\chronologyorientationrow}[2]{%
  \ifcsname triptychchronologyorientation@#1\endcsname
    \csname triptychchronologyorientation@#1\endcsname{#2}%
  \else
    \PackageError{triptych}{No chronology orientation row for #1}{}%
  \fi}
\expandafter\def\csname triptychchronologyorientation@row-4b616c0392b66158\endcsname#1{%
Entrance antiphon: Cf. Psalm 16:6,8; Communion: psalm option: Cf. Psalm 32:18-19 & Historical setting. #1 & Occasion date unresolved\\[.25em]%
}
\expandafter\def\csname triptychchronologyorientation@row-bc236d31cb650062\endcsname#1{%
Entrance antiphon: Cf. Psalm 16:6,8; Responsorial Psalm: Psalm 96:1,3,4-5,7-8,9,10a,10c, Psalm 96:7b; Communion: psalm option: Cf. Psalm 32:18-19 & Composition: Psalms. #1 & \chronologyannotationreach{Ps.16.6}{Ps}{true}\chronologyannotationreach{Ps.16.8}{Ps}{true}\chronologyannotationreach{Ps.32.18}{Ps}{true}\chronologyannotationreach{Ps.32.19}{Ps}{true}\chronologyannotationreach{Ps.95.1}{Ps}{true}\chronologyannotationreach{Ps.95.10}{Ps}{true}\chronologyannotationreach{Ps.95.3}{Ps}{true}\chronologyannotationreach{Ps.95.4}{Ps}{true}\chronologyannotationreach{Ps.95.5}{Ps}{true}\chronologyannotationreach{Ps.95.7}{Ps}{true}\chronologyannotationreach{Ps.95.8}{Ps}{true}\chronologyannotationreach{Ps.95.9}{Ps}{true}\chronologyannotationclaim{critical.psalms.latest-composition-boundary}{composition}{catholic-critical-v1}{preferred}{before the Maccabean period, around 165 B.C.; no individual psalm can be dated securely}{Before c. 165 B.C.}\\[.25em]%
}
\expandafter\def\csname triptychchronologyorientation@row-b4d35a91580c9926\endcsname#1{%
Communion: psalm option: Cf. Psalm 32:18-19 & Traditional attribution: David's reign. #1 & Disputed: \chronologyannotationreach{Ps.32.18}{Ps.32}{true}\chronologyannotationreach{Ps.32.19}{Ps.32}{true}\chronologyannotationclaim{israel.monarchy.david-traditional-era}{traditional-attribution}{catholic-traditional-v1}{disputed}{reigned from 1055 to 1015 B.C.}{B.C. 1055--1015}\\[.25em]%
}
\expandafter\def\csname triptychchronologyorientation@row-84e491e550e8ff0f\endcsname#1{%
Responsorial Psalm: Psalm 96:1,3,4-5,7-8,9,10a,10c, Psalm 96:7b & Superscription setting: The raising of the second Temple by Zorobabel. #1 & \chronologyannotationreach{Ps.95.1}{Ps.95}{true}\chronologyannotationreach{Ps.95.10}{Ps.95}{true}\chronologyannotationreach{Ps.95.3}{Ps.95}{true}\chronologyannotationreach{Ps.95.4}{Ps.95}{true}\chronologyannotationreach{Ps.95.5}{Ps.95}{true}\chronologyannotationreach{Ps.95.7}{Ps.95}{true}\chronologyannotationreach{Ps.95.8}{Ps.95}{true}\chronologyannotationreach{Ps.95.9}{Ps.95}{true}\chronologyannotationclaim{israel.restoration.temple-raised-by-zorobabel}{superscription-setting}{catholic-traditional-v1}{preferred}{(537 B.C.)}{B.C. 537}\\[.25em]%
}
\expandafter\def\csname triptychchronologyorientation@row-f9e687a157cfa9c4\endcsname#1{%
Responsorial Psalm: Psalm 96:1,3,4-5,7-8,9,10a,10c, Psalm 96:7b & Prophetic referent: The Nativity of Our Lord. #1 & Disputed: \chronologyannotationreach{Ps.95.1}{Ps.67 Ps.95 Ps.96 Ps.97}{true}\chronologyannotationreach{Ps.95.10}{Ps.67 Ps.95 Ps.96 Ps.97}{true}\chronologyannotationreach{Ps.95.3}{Ps.67 Ps.95 Ps.96 Ps.97}{true}\chronologyannotationreach{Ps.95.4}{Ps.67 Ps.95 Ps.96 Ps.97}{true}\chronologyannotationreach{Ps.95.5}{Ps.67 Ps.95 Ps.96 Ps.97}{true}\chronologyannotationreach{Ps.95.7}{Ps.67 Ps.95 Ps.96 Ps.97}{true}\chronologyannotationreach{Ps.95.8}{Ps.67 Ps.95 Ps.96 Ps.97}{true}\chronologyannotationreach{Ps.95.9}{Ps.67 Ps.95 Ps.96 Ps.97}{true}\chronologyannotationclaim{life-of-christ.nativity}{prophetic-referent}{catholic-traditional-v1}{disputed}{the year of Rome 750 which he styles 3 B.C.}{B.C. 3}; \chronologyannotationreach{Ps.95.1}{Ps.67 Ps.95 Ps.96 Ps.97}{true}\chronologyannotationreach{Ps.95.10}{Ps.67 Ps.95 Ps.96 Ps.97}{true}\chronologyannotationreach{Ps.95.3}{Ps.67 Ps.95 Ps.96 Ps.97}{true}\chronologyannotationreach{Ps.95.4}{Ps.67 Ps.95 Ps.96 Ps.97}{true}\chronologyannotationreach{Ps.95.5}{Ps.67 Ps.95 Ps.96 Ps.97}{true}\chronologyannotationreach{Ps.95.7}{Ps.67 Ps.95 Ps.96 Ps.97}{true}\chronologyannotationreach{Ps.95.8}{Ps.67 Ps.95 Ps.96 Ps.97}{true}\chronologyannotationreach{Ps.95.9}{Ps.67 Ps.95 Ps.96 Ps.97}{true}\chronologyannotationclaim{life-of-christ.nativity}{prophetic-referent}{catholic-traditional-v1}{disputed}{probably the year 7 B.C.}{c. B.C. 7}\\[.25em]%
}
\expandafter\def\csname triptychchronologyorientation@row-2bb250a9f1f9d49e\endcsname#1{%
First reading: Isaiah 45:1,4-6 & Traditional attribution: Isaias (ministry in Souvay's traditional account). #1 & \chronologyannotationreach{Is.45.1}{Is}{true}\chronologyannotationreach{Is.45.4}{Is}{true}\chronologyannotationreach{Is.45.5}{Is}{true}\chronologyannotationreach{Is.45.6}{Is}{true}\chronologyannotationclaim{israel.prophets.isaias-traditional-era}{traditional-attribution}{catholic-traditional-v1}{preferred}{from the closing year of Ozias, King of Juda (740 B.C.); no prophecies appear to be later than 701 B.C.}{B.C. 740--701}\\[.25em]%
}
\expandafter\def\csname triptychchronologyorientation@row-4ac82ee811fc6241\endcsname#1{%
First reading: Isaiah 45:1,4-6 & Prophecy given: The prophetic activity behind Isaiah 40-55 in the NABRE introduction. #1 & \chronologyannotationreach{Is.45.1}{Is.40 Is.41 Is.42 Is.43 Is.44 Is.45 Is.46 Is.47 Is.48 Is.49 Is.50 Is.51 Is.52 Is.53 Is.54 Is.55}{true}\chronologyannotationreach{Is.45.4}{Is.40 Is.41 Is.42 Is.43 Is.44 Is.45 Is.46 Is.47 Is.48 Is.49 Is.50 Is.51 Is.52 Is.53 Is.54 Is.55}{true}\chronologyannotationreach{Is.45.5}{Is.40 Is.41 Is.42 Is.43 Is.44 Is.45 Is.46 Is.47 Is.48 Is.49 Is.50 Is.51 Is.52 Is.53 Is.54 Is.55}{true}\chronologyannotationreach{Is.45.6}{Is.40 Is.41 Is.42 Is.43 Is.44 Is.45 Is.46 Is.47 Is.48 Is.49 Is.50 Is.51 Is.52 Is.53 Is.54 Is.55}{true}\chronologyannotationclaim{israel.exile.second-isaias-oracles}{prophecy-given}{catholic-critical-v1}{preferred}{toward the end of the Babylonian exile}{Toward the end of the Babylonian exile}\\[.25em]%
}
\expandafter\def\csname triptychchronologyorientation@row-a180ccaa931e671d\endcsname#1{%
Gospel: Matthew 22:15-21 & Event. #1 & Date unresolved\\[.25em]%
}
\expandafter\def\csname triptychchronologyorientation@row-5396102e03d2c0da\endcsname#1{%
Gospel: Matthew 22:15-21 & Composition: The Gospel of St Matthew. #1 & Disputed: \chronologyannotationreach{Matt.22.15}{Matt}{true}\chronologyannotationreach{Matt.22.16}{Matt}{true}\chronologyannotationreach{Matt.22.17}{Matt}{true}\chronologyannotationreach{Matt.22.18}{Matt}{true}\chronologyannotationreach{Matt.22.19}{Matt}{true}\chronologyannotationreach{Matt.22.20}{Matt}{true}\chronologyannotationreach{Matt.22.21}{Matt}{true}\chronologyannotationclaim{composition.gospel-of-matthew}{composition}{catholic-traditional-v1}{disputed}{about A.D. 38-45}{c. A.D. 38--45}; \chronologyannotationreach{Matt.22.15}{Matt}{true}\chronologyannotationreach{Matt.22.16}{Matt}{true}\chronologyannotationreach{Matt.22.17}{Matt}{true}\chronologyannotationreach{Matt.22.18}{Matt}{true}\chronologyannotationreach{Matt.22.19}{Matt}{true}\chronologyannotationreach{Matt.22.20}{Matt}{true}\chronologyannotationreach{Matt.22.21}{Matt}{true}\chronologyannotationclaim{composition.gospel-of-matthew}{composition}{catholic-traditional-v1}{disputed}{about the year 40-42}{c. A.D. 40--42}; \chronologyannotationreach{Matt.22.15}{Matt}{true}\chronologyannotationreach{Matt.22.16}{Matt}{true}\chronologyannotationreach{Matt.22.17}{Matt}{true}\chronologyannotationreach{Matt.22.18}{Matt}{true}\chronologyannotationreach{Matt.22.19}{Matt}{true}\chronologyannotationreach{Matt.22.20}{Matt}{true}\chronologyannotationreach{Matt.22.21}{Matt}{true}\chronologyannotationclaim{composition.gospel-of-matthew}{composition}{catholic-traditional-v1}{disputed}{the years 40-45}{A.D. 40--45}; \chronologyannotationreach{Matt.22.15}{Matt}{true}\chronologyannotationreach{Matt.22.16}{Matt}{true}\chronologyannotationreach{Matt.22.17}{Matt}{true}\chronologyannotationreach{Matt.22.18}{Matt}{true}\chronologyannotationreach{Matt.22.19}{Matt}{true}\chronologyannotationreach{Matt.22.20}{Matt}{true}\chronologyannotationreach{Matt.22.21}{Matt}{true}\chronologyannotationclaim{composition.gospel-of-matthew}{composition}{catholic-traditional-v1}{disputed}{about the year 60-68}{c. A.D. 60--68}; \chronologyannotationreach{Matt.22.15}{Matt}{true}\chronologyannotationreach{Matt.22.16}{Matt}{true}\chronologyannotationreach{Matt.22.17}{Matt}{true}\chronologyannotationreach{Matt.22.18}{Matt}{true}\chronologyannotationreach{Matt.22.19}{Matt}{true}\chronologyannotationreach{Matt.22.20}{Matt}{true}\chronologyannotationreach{Matt.22.21}{Matt}{true}\chronologyannotationclaim{composition.gospel-of-matthew}{composition}{catholic-traditional-v1}{disputed}{about the years 64-67}{c. A.D. 64--67}; \chronologyannotationreach{Matt.22.15}{Matt}{true}\chronologyannotationreach{Matt.22.16}{Matt}{true}\chronologyannotationreach{Matt.22.17}{Matt}{true}\chronologyannotationreach{Matt.22.18}{Matt}{true}\chronologyannotationreach{Matt.22.19}{Matt}{true}\chronologyannotationreach{Matt.22.20}{Matt}{true}\chronologyannotationreach{Matt.22.21}{Matt}{true}\chronologyannotationclaim{composition.gospel-of-matthew}{composition}{catholic-traditional-v1}{disputed}{about the year 50}{c. A.D. 50}\\[.25em]%
}
\expandafter\def\csname triptychchronologyorientation@row-4f74e963ab3022f1\endcsname#1{%
Communion: Mark option: Mark 10:45 & Composition: The Gospel of St Mark. #1 & Disputed: \chronologyannotationreach{Mark.10.45}{Mark}{true}\chronologyannotationclaim{composition.gospel-of-mark}{composition}{catholic-traditional-v1}{disputed}{the tenth or twelfth year after the Ascension (A.D. 38-40)}{A.D. 38--40}; \chronologyannotationreach{Mark.10.45}{Mark}{true}\chronologyannotationclaim{composition.gospel-of-mark}{composition}{catholic-traditional-v1}{disputed}{A.D. 40}{A.D. 40}; \chronologyannotationreach{Mark.10.45}{Mark}{true}\chronologyannotationclaim{composition.gospel-of-mark}{composition}{catholic-traditional-v1}{disputed}{the third year of Claudius (A.D. 43)}{A.D. 43}; \chronologyannotationreach{Mark.10.45}{Mark}{true}\chronologyannotationclaim{composition.gospel-of-mark}{composition}{catholic-traditional-v1}{disputed}{somewhere between A.D. 50 and 67}{A.D. 50--67}; \chronologyannotationreach{Mark.10.45}{Mark}{true}\chronologyannotationclaim{composition.gospel-of-mark}{composition}{catholic-traditional-v1}{disputed}{between 60 and 70}{A.D. 60--70}; \chronologyannotationreach{Mark.10.45}{Mark}{true}\chronologyannotationclaim{composition.gospel-of-mark}{composition}{catholic-traditional-v1}{disputed}{after Matthew, and not as late as the destruction of Jerusalem}{After Matthew, and not as late as the destruction of Jerusalem}\\[.25em]%
}
\expandafter\def\csname triptychchronologyorientation@row-74dca80fab4f2c56\endcsname#1{%
Gospel acclamation: Philippians 2:15d,16a & Composition: The Epistle to the Philippians. #1 & Disputed: \chronologyannotationreach{Phil.2.15}{Phil}{true}\chronologyannotationreach{Phil.2.16}{Phil}{true}\chronologyannotationclaim{composition.epistle-to-the-philippians}{composition}{catholic-traditional-v1}{disputed}{(Philemon; Colossians; Ephesians; Philippians), 61}{A.D. 61}; \chronologyannotationreach{Phil.2.15}{Phil}{true}\chronologyannotationreach{Phil.2.16}{Phil}{true}\chronologyannotationclaim{composition.epistle-to-the-philippians}{composition}{catholic-traditional-v1}{disputed}{at Rome (A.D. 62-64)}{A.D. 62--64}\\[.25em]%
}
\expandafter\def\csname triptychchronologyorientation@row-d67f41bffc530c4a\endcsname#1{%
Second reading: 1 Thessalonians 1:1-5b & Composition: The First Epistle to the Thessalonians. #1 & Disputed: \chronologyannotationreach{1Thess.1.1}{1Thess}{true}\chronologyannotationreach{1Thess.1.2}{1Thess}{true}\chronologyannotationreach{1Thess.1.3}{1Thess}{true}\chronologyannotationreach{1Thess.1.4}{1Thess}{true}\chronologyannotationreach{1Thess.1.5}{1Thess}{true}\chronologyannotationclaim{composition.first-epistle-to-the-thessalonians}{composition}{catholic-traditional-v1}{disputed}{(1 and 2 Thessalonians), 52}{A.D. 52}; \chronologyannotationreach{1Thess.1.1}{1Thess}{true}\chronologyannotationreach{1Thess.1.2}{1Thess}{true}\chronologyannotationreach{1Thess.1.3}{1Thess}{true}\chronologyannotationreach{1Thess.1.4}{1Thess}{true}\chronologyannotationreach{1Thess.1.5}{1Thess}{true}\chronologyannotationclaim{composition.first-epistle-to-the-thessalonians}{composition}{catholic-traditional-v1}{disputed}{in the year 53 or 54 according to the commonly received scheme of Pauline chronology}{A.D. 53--54}\\[.25em]%
}

Composition dates concern the writing of a work. Narrated events, received attributions, superscription settings and Christian prophetic referents answer different questions. The historical Catholic sources below preserve alternative judgments rather than a single agreed chronology. The modern NABRE introductions supply the identified critical horizons for the Psalms and Isaiah. Neither the date of the Sunday nor the date of a Missal is a date for its biblical passages.

\subsection{Psalms of shelter, mercy and universal praise}
\label{history-psalms}
\textbf{Entrance: cf. Psalm 16:6,8 (Vulgate; Hebrew 17).}
\chronodate{entrance}{\chronologyannotation{entrance}}

The Psalm's received title names a prayer of David but identifies no particular occasion. Enemies and sheltering wings do not date a battle, flight or place of composition. The place and occasion of this prayer remain unresolved. Augustine's Christ-and-body speaker is an interpretation of its voice, not a newly dated event in David's biography.

\textbf{Communion option: cf. Psalm 32:18--19 (Vulgate; Hebrew 33).}
\chronodate{communion-psalm}{\chronologyannotation{communion-psalm}}

The Davidic heading in the Douay and Clementine traditions likewise names no event. The reign era in Corbett's \work{King David} is explicitly his ``usual chronology''; his opening account also acknowledges later datings influenced by Assyrian evidence. That disputed reference era dates the attributed author's kingship, not this Psalm's composition. Israel and Judah are its broad geographic frame, not a recovered room in which the poem was written.\footnote{John Corbett, ``King David,'' \work{Catholic Encyclopedia}, IV (1908), opening chronology and account of David.}

For all three Psalms the common latest boundary comes from the NABRE Psalms introduction, which states the uncertainty of dating individual Psalms. It does not assert that they were written together or precisely at the displayed boundary. Their individual writing places are likewise uncertain.\footnote{USCCB, NABRE, ``The Book of Psalms,'' introduction, paragraph on dating and collection; official web delivery checked 13 October 2026.}

\textbf{Responsorial Psalm: Hebrew 96 / Vulgate 95.}
\chronodate{psalm}{\chronologyannotation{psalm}}

The Greek and Vulgate superscription refers to a house built after captivity; the Hebrew title does not carry the same superscription. The historical association above is with Zorobabel's raising of Jerusalem's second Temple, as described at the opening of the encyclopedia's Temple article. The superscription does not securely distinguish that raising from the later completion. Its associated date is therefore not a demonstrated composition date or an unambiguous date of dedication.\footnote{Barnabas Meistermann, ``Temple of Jerusalem,'' \work{Catholic Encyclopedia}, XIV (1912), opening historical account, together with the Psalm's superscription.}

The Nativity alternatives concern the Psalm's Christian prophetic referent. Gerard's \work{General Chronology} reports a correction attributed to interpreters of Dionysius; James Howlett's \work{Biblical Chronology} weighs Herod and infancy evidence and reaches a probable date. Both answers remain disputed. The Psalm itself gives neither a Nativity year nor a locality for the birth. Augustine's Christian reading and John Paul II's account of the cross therefore belong to reception, not to the dating of an ancient sanctuary scene.\footnote{John Gerard, ``General Chronology,'' and James Howlett, ``Biblical Chronology,'' \work{Catholic Encyclopedia}, III (1908), Christian-era and Nativity discussions.}

\subsection{Isaiah's address to Cyrus}
\label{history-isaiah}
\chronodate{first-reading}{\chronologyannotation{first-reading}}

The traditional ministry horizon in Souvay's \work{Isaias} concerns the prophet in Judah and Jerusalem and the attribution of the book to him. It is not a separately established date of an encounter between Isaiah and Cyrus. The NABRE introduction treats the activity behind chapters 40--55 within the late Babylonian exile. That is a different account of the prophetic horizon, not the same event shifted within one scheme. Nor does it imply that the final arrangement of every word of the book was completed at once.\footnote{Charles Souvay, ``Isaias,'' \work{Catholic Encyclopedia}, VIII (1910), life and ministry; USCCB, NABRE, ``The Book of Isaiah,'' introduction, discussion of chapters 40--55.}

Babylonian exile supplies the critical horizon; Jerusalem and its Temple are the destination of the promised restoration in Isaiah 44:26--28. Neither answer above dates Cyrus's capture of Babylon or a decree. No such year is needed to identify the appointed oracle's argument. The distinction between a named ruler and the Lord who calls him survives both source frameworks.

\subsection{Matthew's temple controversy}
\label{history-matthew}
\chronodate{gospel}{\chronologyannotation{gospel}}

Matthew 21:23 and the continuing controversies establish the Jerusalem temple setting of the tribute question. Its exact narrated date remains unresolved here. The composition alternatives date the Gospel, not this exchange. The temple is the scene of the narrative; it does not identify the writer's place.

Jacquier's Matthew article reports competing early testimony and the subsequent dates drawn from it. He notes uncertainty in interpreting Irenaeus's account. The first five alternatives above belong to that discussion, which distinguishes early reported dates from a later proposed interval and another interpretation centred on the preaching and deaths of Peter and Paul. The sixth comes from Durand's \work{The New Testament} article's original-Matthew account, accompanied by uncertainty about the Greek rendering. Overlap between intervals does not make them independent confirmations of one exact year.\footnote{Eugène Jacquier, ``Gospel of St. Matthew,'' \work{Catholic Encyclopedia}, X (1911), ``Date and Place of Composition''; Alfred Durand, ``The New Testament,'' XIV (1912), origins of the Gospels.}

\subsection{Mark's journey toward Jerusalem}
\label{history-mark}
\chronodate{communion-mark}{\chronologyannotation{communion-mark}}

Mark 10:32--45 places the teaching on service and ransom on the ascent toward Jerusalem, after the Passion prediction and the request for seats of honour. Verse 46 reaches Jericho after the saying. The precise roadside locality of the teaching is not named. None of the composition dates above supplies an event date for that journey or for the Passion.

MacRory's Mark article reports early subscriptions and chronicles, considers their possible derivation and weighs conflicting testimony toward a broader preferred interval. Its reported early dates and its preferred interval remain distinguishable answers. Durand's \work{The New Testament} article gives the other numerical interval. The Pontifical Biblical Commission's responsa V--VI of 26 June 1912 supply relative order and a limit before Jerusalem's destruction or the beginning of its siege, not a numerical year. A relative answer must remain relative.\footnote{Joseph MacRory, ``Gospel of Saint Mark,'' \work{Catholic Encyclopedia}, IX (1910), ``Time of Composition''; Durand, ``The New Testament''; PBC, responsa on Mark and Luke, V--VI, \work{Acta Apostolicae Sedis} 4 (1912), pp.463--465.}

\subsection{Philippians: writer and recipients}
\label{history-philippians}
\chronodate{acclamation}{\chronologyannotation{acclamation}}

The two historical composition schemes give different dates within a Roman-captivity account. Prat's schedule supplies the first; the encyclopedia's separate Philippians article supplies the second and argues for Rome rather than Caesarea, while discussing the stage of captivity. Rome is the writer's location in that account, Philippi the recipients' city. The account is not a claim that every scholarly proposal places the writing at Rome. The short acclamation does not date its own delivery as an exhortation.\footnote{Ferdinand Prat, ``St. Paul,'' \work{Catholic Encyclopedia}, XI (1911), chronology; Achille Vander Heeren, ``Epistle to the Philippians,'' XII (1911), occasion and date/place of composition.}

\subsection{Thessalonians: thanksgiving sent from Corinth}
\label{history-thessalonians}
\chronodate{second-reading}{\chronologyannotation{second-reading}}

Prat's schedule supplies one traditional scheme; Drum's article supplies the other. Drum argues for Corinth by joining the letter's three senders, Acts 18:5 and the report brought by Timothy. Thessalonica is the receiving community's city. The writing date does not automatically date every act of faith and love remembered in the thanksgiving. Drum also discusses a critical alternative outside the returned comparison; it is not silently made a third date here.\footnote{Prat, ``St. Paul,'' chronology; Walter Drum, ``Epistles to the Thessalonians,'' \work{Catholic Encyclopedia}, XIV (1912), time, place and occasion of writing.}

The Collect, Prayer over the Offerings and Prayer after Communion are composed prayers with no appointed biblical citation. They receive no biblical chronology from a resemblance to a verse. The Scripture references of the Missal antiphons identify adaptations or an excerpt; their dates concern the underlying passages, not the composition of those liturgical adaptations.
